% TOP_PROBLEM: {"id": 30001574, "title": "Morrison’s Cone Conjecture for Calabi–Yau Pairs", "category_id": 5, "category": {"id": 5, "name": "algebraic_geometry", "display_name": "Algebraic Geometry", "description": "Geometric objects defined by polynomial equations.", "slug": "algebraic-geometry", "order_index": 5, "created_at": "2026-07-31T15:26:25.671Z"}, "status": "partially_solved"}
% FIRST_POSTED: 2026-10-01
% ATTEMPT_STATUS: unresolved
% !TeX program = lualatex
% Definition and sources only; this file is not an LLM solution attempt.
\documentclass[11pt]{article}
\usepackage[margin=25mm]{geometry}
\usepackage{fontspec}
\setmainfont{Latin Modern Roman}
\usepackage{amsmath,amssymb,mathtools}
\usepackage{unicode-math}
\setmathfont{Latin Modern Math}
\usepackage{xurl,hyperref}
\hypersetup{colorlinks=true,urlcolor=blue}
\setcounter{secnumdepth}{0}
\setlength{\parindent}{0pt}
\setlength{\parskip}{5pt}
\setlength{\emergencystretch}{3em}
\begin{document}
\section*{\texorpdfstring{Morrison’s Cone Conjecture for Calabi–Yau Pairs}{Morrison’s Cone Conjecture for Calabi–Yau Pairs}}\label{30001574}
\subsection{Definitions and mathematical statement}
Let $(X,\Delta)$ be a projective complex ${\mathbb{Q}}$-factorial klt pair with $K_X+\Delta\equiv0$. Here ${\mathbb{Q}}$-factorial means every Weil divisor has a Cartier multiple, and klt means discrepancies exceed $-1$. In $N^1(X)_{{\mathbb{R}}}$, let $\operatorname{Nef}^e$ be the effective nef cone and $\operatorname{Mov}^e$ the effective movable cone in the cited convention. Automorphisms preserve the pair; pseudo-automorphisms are birational self-maps preserving it and isomorphic in codimension one. The two assertions ask that each relevant cone be covered by translates of a rational polyhedral cone under its indicated group, with interiors of distinct translates disjoint in the standard fundamental-domain sense. Rational polyhedral means generated by finitely many rational divisor classes. Both the nef and movable assertions are part of the target.
\par\smallskip\textit{Formulation and concept references:} \href{https://arxiv.org/abs/2406.07307}{[S1]}.

\subsection{Short English statement}
Do automorphisms and pseudo-automorphisms reduce the effective nef and movable cones of every klt Calabi–Yau pair to rational polyhedral fundamental domains?
\subsection{Research attempt: Picard rank one and a conditional rank-two fundamental domain}
\textbf{Status.} We prove both requested cone assertions in Picard rank one and an explicit rank-two tiling lemma. The latter includes hypotheses on the actual cone and group action; it does not derive them from the Calabi--Yau assumptions. The primary source [S1] studies relations between cone-conjecture formulations and examples, rather than giving an unconditional result for every pair.

\subsubsection{1. Both cones when the Picard number is one}
Let $\rho(X)=1$ and choose an ample Cartier class $H$. The numerical space is $\mathbb RH$. Any nonzero effective divisor class $aH$ has $a>0$: its intersection with $H^{\dim X-1}$ is positive, whereas $H^{\dim X}>0$. The ample class itself is effective after multiplication by a positive integer. Thus the effective nef cone and the effective movable cone, in either usual effective-cone intersection or effective-generator convention, are both
\[\mathbb R_{\ge0}H.\]
For movability, a sufficiently large multiple of $H$ is base-point free, so the entire ray is included; no negative ray is effective. An automorphism or pseudo-automorphism acts integrally and invertibly on the divisor lattice modulo numerical equivalence and torsion. A rank-one lattice has only the two automorphisms $+1,-1$, and preservation of the effective ray forces $+1$. Therefore the whole ray is a rational polyhedral fundamental domain for either image group. A kernel acting trivially on classes does not produce distinct translates.

\subsubsection{2. A rank-two tiling calculation}
Let $V$ be a rational two-dimensional vector space and $g\in\mathrm{GL}(V_{\mathbb Z})$ have positive real eigenvalues $\lambda>1$ and $\lambda^{-1}$, with eigenvectors $v_+,v_-$. Suppose an invariant cone has interior
\[C^\circ=\{xv_++yv_-:x,y>0\}.\]
Choose a rational interior vector $D=av_++bv_-$, $a,b>0$, and put $r_0=a/b$. The rays through $g^nD$ have ratios $r_n=\lambda^{2n}r_0$. For every $r>0$ there is an integer $n$ with $r_n\le r\le r_{n+1}$. Therefore
\[C^\circ\cup\{0\}=\bigcup_{n\in\mathbb Z}g^n\operatorname{Cone}(D,gD).\tag{1}\]
The interiors of adjacent sectors have disjoint ratio intervals, and all other sectors are separated. Each sector is rational polyhedral because both $D$ and $gD$ are rational. This proves an explicit fundamental-domain statement for the cyclic group $\langle g\rangle$ on the open cone plus its origin.

\subsubsection{3. Boundary and group hypotheses cannot be omitted}
The boundary eigenrays have ratio zero or infinity and never enter a finite sector in (1). If the eigenrays are irrational, neither contains a nonzero rational class. A cone generated by effective integral divisors can then contain no nonzero boundary element: any finite positive sum on an extremal boundary ray would force every summand onto that same ray. Under this additional effective-cone description, the cone in question is exactly its interior plus zero, and (1) covers it. A closed nef cone containing irrational boundary rays is a different set; the same finite-sector union does not cover those rays.

Likewise, having one automorphism $g$ does not identify the full automorphism or pseudo-automorphism group. If the image group is exactly $\langle g\rangle$, (1) is the desired fundamental domain. If there are additional elements, their translates can overlap sector interiors. One must analyze that larger group or subdivide a domain; merely finding an infinite cyclic subgroup does not prove the full group's assertion. As an opposite obstruction to this approach, a trivial group cannot have a rational polyhedral fundamental domain equal to a genuinely nonpolyhedral cone.

\subsubsection{4. Verification and exact missing geometry}
The ratio transformation is $x/y\mapsto\lambda^2x/y$, so powers range in both directions and cover every finite positive ratio. Rationality is required of $D$, not of the eigenvectors. These checks prove the tiling lemma with its stated assumptions. For higher Picard number, and even for general rank two, the geometry must still supply the shape of each effective cone and enough pair-preserving automorphisms or pseudo-automorphisms. No argument here supplies that information for an arbitrary klt Calabi--Yau pair.

\subsection{Sources}
\begin{itemize}
\item[S1]The effective cone conjecture for Calabi-Yau pairs.
\url{https://arxiv.org/abs/2406.07307}.
\textit{Formulation source.}
\end{itemize}
\end{document}
